%%%%%%%%%%  scan idx 362  |  folio 349  %%%%%%%%%%

On
peut donc de façon évidente étendre \(\ch\) en un homomorphisme de \(K\)-algèbres : \(K\times(1+\widehat A^+)\longrightarrow K\oplus\widehat A_{\mathbb Q}^{+}\). On a donc :
\begin{equation}\tag{6.3.3}
  \ch(a,\varphi)=a+\eta(\Log\varphi).
\end{equation}

\textbf{Définition 6.4.} L'homomorphisme \(\ch:\Ch A\longrightarrow K\oplus\widehat A_{\mathbb Q}^{+}\) qui vient d'être défini s'appelle \emph{caractère de Chern}.

Si \(A\) est de caractéristique \(0\) (i.e. est une \(\mathbb Q\)-algèbre), le caractère de Chern est un isomorphisme~: en effet, soit \(\psi\in K\oplus\widehat A^+\)~; on peut écrire \(\psi=a+\psi_1\), où \(\psi_1\) est sans terme de degré nul~; on définit alors un morphisme inverse de \(\ch\) par \(\psi\rightsquigarrow(a,\exp(\eta^{-1}(\psi_1)))\).

\textbf{Proposition 6.5.} Soit \(A\) une \(K\)-algèbre graduée~; considérons sur \(\Ch(A)\) la filtration définie comme suit (filtration par le degré) :
\begin{equation}\tag{6.5.1}
  (\Ch(A))_0=\Ch(A)~;\qquad
  \forall n\ge1,\ (\Ch(A))_n=0\times 1+\widehat A^{\ge n},
\end{equation}
où \(\widehat A^{\ge n}=\prod_{i\ge n}A^i\), On a les congruences suivantes :

\begin{enumerate}[label=\roman*)]
\item Soient \(n\ge1\), et \(\varphi\in(\Ch A)_n\)~; posons \(\varphi=1+a_n+\cdots\)~; alors :
\begin{equation}\tag{6.5.2}
  \ch(\varphi)=\frac{(-1)^{n-1}}{(n-1)!}.a_n+..~;
\end{equation}

\item Soient \(\varphi\in(\Ch A)_m\), \(\psi\in(\Ch A)_n\)~; posons
\(\varphi=1+a_m+\cdots\), \(\psi=1+b_n+\cdots\)~; alors :
\begin{equation}\tag{6.5.3}
  \varphi*\psi
  =
  1-\frac{(m+n-1)!}{(m-1)!(n-1)!}.a_m b_n+..,
\end{equation}
et par suite, la filtration par le degré est une filtration d'anneau.
\end{enumerate}

i) Le caractère de Chern est par construction un morphisme de foncteurs filtrés~; de la congruence \(\varphi\equiv1+a_n\ \mathrm{mod.}\ (\Ch A)_{n+1}\), on déduit
